% TeXLeaf 可编辑中文 Beamer 模板。
\documentclass[UTF8,aspectratio=169]{ctexbeamer}

\usetheme{CambridgeUS}
\usecolortheme{whale}

\usepackage{mathtools,amssymb}
\usepackage{booktabs}
\usepackage{graphicx}

% Beamer 已提供 theorem、corollary、lemma、definition 和 example。
\theoremstyle{plain}
\newtheorem{axiom}[theorem]{公理}
\newtheorem{proposition}[theorem]{命题}
\newtheorem{claim}[theorem]{断言}
\newtheorem{conjecture}[theorem]{猜想}
\newtheorem{hypothesis}[theorem]{假设}

\theoremstyle{definition}
\newtheorem{assumption}[theorem]{假定}

\theoremstyle{example}
\newtheorem{exercise}[theorem]{练习}

\theoremstyle{remark}
\newtheorem{remark}[theorem]{注}

\title{@{0:演示文稿题目}}
\subtitle{@{1:副标题（可选）}}
\author{@{2:作者姓名}}
\institute{@{3:单位}}
\date{\today}

\AtBeginSection[]{
  \begin{frame}{目录}
    \tableofcontents[currentsection]
  \end{frame}
}

\begin{document}

\begin{frame}
  \titlepage
\end{frame}

\begin{frame}{目录}
  \tableofcontents
\end{frame}

\section{引言}

\begin{frame}{@{4:研究动机}}
  \begin{itemize}
    \item @{5:第一点}
    \item @{6:第二点}
  \end{itemize}
\end{frame}

\section{主要结果}

\begin{frame}{@{7:主要结果}}
  @{8:在此加入核心论证、定理、图片或表格。}
\end{frame}

\section{结论}

\begin{frame}{结论}
  @{9:总结主要内容与后续工作。}
\end{frame}

@10
\end{document}
